\metatitle{Symplectic and orthogonal Schur functions}
\metadescription{Symplectic and orthogonal Schur polynomials, their tableau
models, determinantal identities, crystals, and branching rules.}
\metakeywords{symplectic Schur polynomial, orthogonal Schur polynomial,
symplectic tableau, orthogonal tableau, Jacobi-Trudi identity}

\section[schurSymplectic]{Symplectic Schur polynomials}


\begin{polydata}{schurSymplectic}
  Name     & Symplectic Schur polynomials \\
  Space    &  Other   \\
  Basis    &  False   \\
  Rating   &  1      \\
  Bib      &  King1971 \\
  Year     &  1971? \\
  Keywords & jacobi-trudi,tableaux,cauchy-identity,skew \\
  Symbol   & $\schurSp_{\lambda/\mu}(\xvec)$ \\
  Category & Schur \\
  TransformsTo & schurOrthogonal | Hamel1997 | map=omega and conjugate index \\
\end{polydata}


The \emph{even symplectic Schur polynomials} were explicitly described
combinatorially in \cite{King1976}, and are characters for the Lie group
$\mathrm{Sp}(2n)$.  See also \cite{Krattenthaler1998}.
\name{Robert Proctor} gives Young-tableau and Gelfand-pattern models, together
with branching rules, for characters of the classical groups \cite{Proctor1994}.
\name{Markus Fulmek} and \name{Christian Krattenthaler} give nonintersecting
lattice-path proofs of Jacobi--Trudi and Giambelli-type identities for
symplectic and orthogonal characters \cite{FulmekKrattenthaler1997}.
\name{Peter S. Campbell} and \name{Anna Stokke} prove hook-content formulae
for symplectic and orthogonal tableaux \cite{CampbellStokke2012}.  These
formulae give product evaluations for the corresponding tableau generating
functions.
\name{Soichi Okada} studies intermediate symplectic characters and
\hyperref[planePartitions]{shifted plane partitions} of shifted double
staircase shape
\cite{Okada2020Intermediate}.  These characters interpolate between Schur
polynomials and symplectic characters.
\name{George Balla} constructs PBW tableaux for symplectic PBW degenerate
\hyperref[flagVarieties]{flag varieties} and gives defining equations for
these varieties
\cite{Balla2022}.

We work in the alphabet
\[
 \xvec = (x_1,x_1^{-1},x_2,x_2^{-1},\dotsc,x_n,x_n^{-1}).
\]
A \defin{symplectic tableau}\label{symplecticTableau} (also known as King
tableau) of shape $\lambda/\mu$ is a filling of the shape $\lambda/\mu$
with entries in
\[
1 \lt \overline{1} \lt 2 \lt \overline 2 \lt \dotsb \lt n \lt \overline{n}
\]
such that entries are weakly increasing along rows and strictly increasing
along columns, and entries in row $i$ are greater than or equal to $i$.
Let $SP(\lambda/\mu)$ denote the set of such tableaux.
Define the \defin{even symplectic Schur polynomials}
\label{evenSymplecticSchurPolynomial}
(using the definition in \cite{Hamel1997}) as
\[
\schurSp_{\lambda/\mu}(\xvec) = \sum_{T \in SP(\lambda/\mu)} \xvec^{T},
\]
where entries $i$ in $T$ contribute with $x_i$, and entries $\overline{j}$
contribute with $x_j^{-1}$.

\begin{example}
The following tableau is an element of $SP(\lambda/\mu)$.
\begin{figure}
\begin{ytableau}
1 & 1 & \overline{1} & \overline{1} & \overline{2} & 3 \\
\overline{1} & \overline{2} & 3 & 4 \\
3 & 3
\end{ytableau}
\end{figure}
\end{example}


For symmetry, see \cite[Thm.~6.12]{SundaramThesis}, where a type of
Bender--Knuth involution is defined.  The proof seems to have a gap; see the
discussion in \cite{MO362997}.  A complete proof is available in
\cite{Gutierrez2024}.


There are also alternative tableau descriptions.  For example, the
\defin{Kashiwara--Nakashima tableaux}\label{kashiwaraNakashimaTableau} carry a
\hyperref[kashiwaraCrystals]{$U_q(\mathfrak{sp}(2m,\setC))$-crystal
structure}; see \cite{AzenhasFellerTorres2022}.


\subsection[schurSymplecticBialternant]{Bialternant formula}

The \emph{even symplectic Schur functions} can be expressed as
\[
\schurSp_\lambda(x^{\pm 1}_{1},x^{\pm 1}_{2},\dotsc, x^{\pm 1}_{n})
=
\frac{
\left| x_i^{\lambda_j+n-j+1} - x_i^{-(\lambda_j+n-j+1)}\right|_{1 \leq i,j \leq n}
}{
\left| x_i^{n-j+1} - x_i^{-(n-j+1)}\right|_{1 \leq i,j \leq n}
}
\]

The \emph{odd symplectic Schur functions} have a more complicated
bialternant formula; see \cite{Okada2020OddSymplectic}.

\name{Naihuan Jing}, \name{Zhijun Li}, \name{Xinyu Pan},
\name{Danxia Wang}, and \name{Chang Ye} introduce skew odd orthogonal
characters and interpolating Schur polynomials \cite{JingLiPanWangYe2025}.
They prove Cauchy and Jacobi--Trudi identities, describe
Gelfand--Tsetlin-pattern models, and give transition formulas among
symplectic, even orthogonal, and odd orthogonal characters.


\subsection[schurSymplecticCauchy]{Cauchy identity}


The following Cauchy identity is proved by \name{Sheila Sundaram} in her
thesis; see \cite{SundaramThesis}.
\[
\sum_{\lambda} \schurSp_\lambda(x^{\pm 1}_{1},x^{\pm 1}_{2},\dotsc, x^{\pm 1}_{n}) \schurS_\lambda(y_1,\dotsc,y_n)
= \prod_{1\leq i \lt j \leq n} (1-y_i y_j)
\prod_{1\leq i \lt j \leq n} \frac{1}{(1-y_i x_j)(1-y_i x^{-1}_j)}.
\]

\name[Ronald C. King]{Ronald C.~King} derives multiplicative generating
functions for series of characters of classical Lie groups
\cite{King2023x}.  The identities generalize Schur and symplectic Cauchy-type
expansions to all classical types, with recurrence relations for the expansion
coefficients coming from Weyl group actions.

\name{Aalekh Patel}, \name{Harsh Patel}, and \name{Anna Stokke} prove
orthosymplectic analogues of the Cauchy and dual Cauchy identities
\cite{PatelPatelStokke2022}.  Their proofs are bijective and use insertion
algorithms modeled on \name{Allan Berele}'s symplectic insertion.

\name{Anna Stokke} gives a Pieri rule for multiplying an orthosymplectic
character by an orthosymplectic character indexed by a one-row partition
\cite{Stokke2018}.  The resulting product rule agrees with
\name{Sheila Sundaram}'s Pieri rule for symplectic characters.

\name{Dan Betea} studies symplectic and orthogonal analogues of
\name{Andrei Okounkov}'s Schur measure \cite{Betea2018x}.  He proves that
these measures have determinantal correlation functions with double-contour
kernels, and relates them to Toeplitz-plus-Hankel determinants and asymptotic
edge kernels.





\subsection[schurSymplecticJacobiTrudi]{Jacobi--Trudi}

In \cite{FulmekKrattenthaler1997}, the following analogues of the
Jacobi--Trudi identities are proved:
\begin{align}
\schurSp_\lambda(x^{\pm 1}_{1},x^{\pm 1}_{2},\dotsc, x^{\pm 1}_{n})
& = \frac{1}{2}
\det[ \completeH_{\lambda_i - i+j}(\xvec)
      + \completeH_{\lambda_i - i- j + 2}(\xvec) ]_{1\leq i,j \leq \length(\lambda)} \\
& = \det[ \elementaryE_{\lambda'_j - j+i}(\xvec) - \elementaryE_{\lambda'_j - j- i}(\xvec) ]_{1\leq i,j \leq \lambda_1}.
\end{align}
These formulas allow us to consider $\schurSp_\lambda(\xvec)$ as a symmetric
function, by instead using the infinite alphabet $x_1,x_2,\dotsc$.

\begin{example}
For example, the monomial expansions begin as follows:
\begin{align}
\schurSp_{4}  &= \monomial_{4} + \monomial_{22} + \monomial_{31}+\monomial_{211}+\monomial_{1111} \\
\schurSp_{31} &= -\monomial_{2} - \monomial_{11} + \monomial_{22}+\monomial_{31}+2\monomial_{211}+3\monomial_{1111} \\
\schurSp_{22} &= - \monomial_{11} + \monomial_{22} + \monomial_{211}+2\monomial_{1111} \\
\schurSp_{211} &= 1 - \monomial_{2} -2 \monomial_{11} + \monomial_{211} + 3\monomial_{1111} \\
\schurSp_{1111} &= - \monomial_{11} + \monomial_{1111}
\end{align}
The finite symplectic alphabet must be substituted in these expressions in
order for these formulas to agree with the tableau definition.
For example,
\begin{align}
\schurSp_{11}(\xvec) &= - 1 + \monomial_{11}(\xvec)\\
\schurSp_{11}(x_1,x_1^{-1},x_2,x_2^{-1}) &=
1 + x_1^{-1}x_2^{-1} + x_1x^{-1}_2+x^{-1}_1x_2 + x_1x_2
\end{align}
\end{example}


\subsection[schurSymplecticGiambelli]{Giambelli formula}

In Frobenius coordinates,
\[
\schurSp_{(\alpha|\beta)} = \left| \schurSp_{(\alpha_i|\beta_j)} \right|_{s \times s}
\]
see, for example, \cite[Eq.~(3.11)]{FulmekKrattenthaler1997}.



\subsection[schurSymplecticCrystal]{Kashiwara crystals}

In \cite{Lee2019}, \name{Seung Jin Lee} defines a
\hyperref[crystalGraphs]{crystal graph} structure on King tableaux,
so that connected components correspond to single symplectic Schur functions.
The crystal structure on $K(\mu,n)$, the set of symplectic tableaux of shape
$\mu$ and maximal entry at most $n$, is isomorphic to the crystal $B(\mu)$,
the irreducible representation of $\mathfrak{sp}_{2n}$ indexed by $\mu$.

Lee also gives a crystal-preserving bijection between symplectic tableaux and
\defin{oscillating tableaux}\label{oscillatingTableau}.

An explicit conjecture regarding the expansion
\[
 \schurSp_\lambda(\xvec^{\pm}) \schurS_\mu(\xvec^{\pm}) = \sum_{\nu} c^{\nu}_{\lambda\mu} \schurSp_\nu(\xvec^{\pm})
\]
appears in \cite[Conj.~6.2]{Lee2019}.


\subsection[schurSymplecticMurnaghanNakayama]{Murnaghan--Nakayama rules}

\name{Nishu Kumari} and \name{Anna Stokke} prove
\hyperref[murnaghanNakayamaRule]{Murnaghan--Nakayama rules} for symplectic,
orthogonal, and orthosymplectic Schur functions \cite{KumariStokke2026}.
For symplectic Schur functions, the product of a symplectic power sum with
$\schurSp_\mu$ is expressed as a sum of three kinds of terms: symplectic Schur
functions obtained by adding a \hyperref[skewShapes]{border strip} to $\mu$,
those obtained by removing a \hyperref[skewShapes]{border strip} from $\mu$,
and an additional Weyl-denominator
correction term.  The orthogonal rules are analogous, while the
orthosymplectic rule has a third term involving both symplectic and ordinary
Schur functions.


\subsection[schurMiscRootTwists]{Root-of-unity twists}

\name{Dudley Littlewood}'s factorization of Schur functions at variables
twisted by a primitive root of unity can be expressed in terms of
\hyperref[partitionQuotient]{cores and quotients}.  \name{Seamus P. Albion}
lifts these factorizations to universal characters
\cite{Albion2023}.  This also gives uniform proofs of related factorizations
for orthogonal and symplectic characters.

\name{Nishu Kumari} extends character factorizations at root-of-unity twists
for general linear, symplectic, orthogonal, and special orthogonal groups
\cite{Kumari2024}.  The paper characterizes when the specialized characters
vanish and factors the nonzero cases into characters of smaller classical
groups, using Weyl character formulas and beta-sets of $t$-core partitions.
\name{Seamus P. Albion} embeds the Schur, symplectic, and orthogonal
factorization formulas into a family involving a parameter $z$ and a
parameter $q$ \cite{Albion2025x}.  The proof uses $t$-cores and $t$-quotients
of $z$-asymmetric partitions and relates the formulas to \hyperref[plethysm]{plethysm}.
\name{Arvind Ayyer} and \name{Nishu Kumari} obtain further root-of-unity
factorizations for classical and universal characters
\cite{AyyerKumari2024x}, including specializations of Schur polynomials and
staircase-shape factorizations.


\subsection[schurSymplecticOtherProperties]{Further properties}

There are analogues of \hyperref[gtpatterns]{Gelfand--Tsetlin patterns} for
these functions; see, for example,
\cite{Proctor1994,AyyerFischer2019,JingLiWang2024}.

\name{Tewodros Amdeberhan}, \name{George E. Andrews}, and
\name{Cristina Ballantine} study partition identities involving
\hyperref[partitionCores]{hook lengths}
and symplectic or orthogonal contents \cite{AmdeberhanAndrewsBallantine2023}.
They prove special cases of Nekrasov--Okounkov-type conjectures and relate
nonzero symplectic contents to partitions into distinct even parts.

\name{Naihuan Jing}, \name{Zhijun Li}, and \name{Danxia Wang} use vertex
operators to realize the skew symplectic and skew orthogonal Schur functions
of \name{Kazuhiko Koike} and \name{Itaru Terada} \cite{JingLiWang2024}.  They
derive Jacobi--Trudi identities, Gelfand--Tsetlin pattern formulas, branching
rules, and Cauchy-type identities for these functions.
\name[V. Sathish Kumar]{V.~Sathish Kumar} and \name{Jacinta Torres} prove a
bijection between the branching models of Kwon and Sundaram using flagged
hives \cite{KumarTorres2024x}.  Their construction uses a hive-model symmetry
for Littlewood--Richardson coefficients and gives a new flagged-hive
branching model.

\name{Olga Azenhas} gives a combinatorial inverse to Watanabe's type-$AII$
quantum Littlewood--Richardson map \cite{Azenhas2026x}.  Using reverse
Schensted column insertion and inverse reduction of symplectic columns, this
proves the $GL_{2n}(\setC)$-to-$Sp_{2n}(\setC)$ branching bijection and
realizes the orthogonal-transpose symmetry on Littlewood--Richardson--Sundaram
tableaux.

\name{Zhihong Jin}, \name{Naihuan Jing}, \name{Zhijun Li}, and
\name{Danxia Wang} introduce universal symplectic and orthogonal functions
which include symplectic Schur functions, odd symplectic characters, and
universal symplectic and orthogonal characters as special cases
\cite{JinJingLiWang2022x}.  Their vertex-operator construction gives skew
versions, transition formulas, and general branching rules.

\name{Nishu Kumari} proves a determinantal formula for orthosymplectic Schur
functions analogous to the Moens--Van der Jeugt determinant for general
linear Lie superalgebras \cite{Kumari2024x}.  As a consequence, the odd
symplectic characters of \name{Robert Proctor} are realized as specialized
orthosymplectic characters.

The even symplectic Schur functions can be expressed as a sum of skew Schur
functions:
\[
\schurSp_\lambda(\xvec) = \sum_\alpha (-1)^{|\alpha|/2} \schurS_{\lambda/\alpha}(\xvec)
\]
where $\alpha$ ranges over partitions whose
\hyperref[prelimPartitions]{Frobenius coordinates} are
$(a_1,a_2,\dotsc, | a_1+1,a_2+1,\dotsc)$,
including the empty partition.


The following branching rule is due to \name{Dudley Littlewood}
\cite[p.~295]{Littlewood1977}; see also \cite[p.~54]{Krattenthaler1998}.
\begin{theorem}[Littlewood]
\[
\schurS_{\lambda}(\xvec)
= \sum_{\nu} \schurSp_\nu(\xvec)
  \sum_{\substack{\mu\\ \mu' \text{ even}}} c^{\lambda}_{\mu,\nu}
\]
where the sum is over all partitions $\mu$ where the columns are even.
\end{theorem}





\section[schurOrthogonal]{Orthogonal Schur polynomials}




\begin{polydata}{schurOrthogonal}
  Name   & Orthogonal Schur polynomials \\
  Space    & Other \\
  Basis    & False \\
  Rating   & 1 \\
  Bib      & King1971 \\
  Year     & 1971? \\
  Keywords & jacobi-trudi,tableaux,cauchy-identity,skew \\
  Symbol   & $\schurOr_\lambda(\xvec)$ \\
  Category & Schur \\
  TransformsTo & schurSymplectic | Hamel1997 | map=omega and conjugate index \\
\end{polydata}


A \defin{special orthogonal tableau}\label{specialOrthogonalTableau} of shape
$\lambda/\mu$ is a filling of the shape $\lambda/\mu$
with entries in
\[
1 \lt \overline{1} \lt 2 \lt \overline 2 \lt \dotsb \lt n \lt \overline{n} \lt \infty
\]
such that (a) entries are weakly increasing along rows
and all finite entries are strictly increasing along columns,
(b) entries in row $i$ are greater than or equal to $i$,
and (c) there is at most one $\infty$ in every row.
Let $SO(\lambda/\mu)$ denote the set of such tableaux.

Define the \defin{even orthogonal Schur polynomials}
\label{evenOrthogonalSchurPolynomial} using the definition in
\cite{Hamel1997}:
\[
\schurOr_{\lambda/\mu}(\xvec) = \sum_{T \in SO(\lambda/\mu)} \xvec^{T},
\]
where entries $i$ in $T$ contribute with $x_i$ and entries $\overline{j}$
contribute with $x_j^{-1}$.
The $\infty$-elements do not contribute.

The \hyperref[standardInvolution]{standard involution on symmetric functions}
satisfies $\omega(\schurSp_\lambda)=\schurOr_{\lambda'}$.


\subsection[schurOrthogonalJacobiTrudi]{Jacobi--Trudi}

The orthogonal Schur polynomials satisfy the following Jacobi--Trudi
identities; see \cite{FulmekKrattenthaler1997,Hamel1997}.
They follow from the symplectic identities using
$\omega(\schurSp_\lambda)=\schurOr_{\lambda'}$.

\begin{align}
\schurOr_\lambda(x^{\pm 1}_{1},x^{\pm 1}_{2},\dotsc, x^{\pm 1}_{n})
&
=\det[ \completeH_{\lambda_i - i+j}(\xvec) - \completeH_{\lambda_i - i- j }(\xvec) ]_{1\leq i,j \leq \length(\lambda)} \\
& =\frac{1}{2}\det[ \elementaryE_{\lambda'_i - i+j}(\xvec) +
\elementaryE_{\lambda'_i - i- j + 2}(\xvec) ]_{1\leq i,j \leq \lambda_1}.
\end{align}
Note that the functions on the right-hand side are also evaluated in
the alphabet $\xvec = (x^{\pm 1}_{1},x^{\pm 1}_{2},\dotsc, x^{\pm 1}_{n})$.


More general identities, analogous to the
\hyperref[schurStripDecompFormula]{Hamel--Goulden identities}, see
\cite{Hamel1997}.



As for the symplectic characters, one has the following expansion.
\begin{theorem}[Littlewood]
\[
\schurS_{\lambda}(\xvec)
= \sum_{\nu} \schurOr_\nu(\xvec)
  \sum_{\substack{\mu\\ \mu \text{ even}}} c^{\lambda}_{\mu,\nu}
\]
where the sum is over all partitions $\mu$ where the rows are even.
\end{theorem}


\name{Naihuan Jing}, \name{Zhijun Li}, and \name{Danxia Wang} give a modern
vertex-operator treatment of skew symplectic and orthogonal Schur functions
\cite{JingLiWang2024}.  Their formulas include Jacobi--Trudi identities,
Gelfand--Tsetlin-pattern models, branching rules, and Cauchy-type identities.
%
% \subsection[schurOrthogonalBialternant]{Bialternant formula}
%
% We have that the \emph{orthogonal Schur functions} can be expressed as
% \[
% \schurOr_\lambda(x^{\pm 1}_{1},x^{\pm 1}_{2},\dotsc, x^{\pm 1}_{n})
% =
% \frac{
% \left| x_i^{\lambda_j+n-j+1} - x_i^{-(\lambda_j+n-j+1)}\right|_{1 \leq i,j \leq n}
% }{
% \left| x_i^{n-j+1} - x_i^{-(n-j+1)}\right|_{1 \leq i,j \leq n}
% }
% \]



\subsection[schurOrthogonalCauchy]{Cauchy identity}

\[
\sum_{\lambda} \schurOr_\lambda(x^{\pm 1}_{1},x^{\pm 1}_{2},\dotsc, x^{\pm 1}_{n}) \schurS_\lambda(y_1,\dotsc,y_n)
=
\prod_{1\leq i \leq j \leq n} (1-y_i y_j)
\prod_{1\leq i \lt j \leq n} \frac{1}{(1-y_i x_j)(1-y_i x^{-1}_j)}.
\]

\subsection[schurOrthogonalOtherProperties]{Further properties}

The \defin{orthogonal Schur functions}\label{orthogonalSchurFunction} can be
expressed as a sum of skew Schur functions as
\[
\schurOr_\lambda(\xvec) = \sum_\beta (-1)^{|\beta|/2} \schurS_{\lambda/\beta}(\xvec)
\]
where $\beta$ ranges over the Frobenius coordinates $(b_1+1,b_2+1,\dotsc, | b_1,b_2,\dotsc)$,
including the empty partition.
